\documentclass[11pt,a4paper]{article}
\usepackage[margin=1in]{geometry}
\usepackage[T1]{fontenc}
\usepackage[utf8]{inputenc}
\usepackage{lmodern}
\usepackage{hyperref}
\usepackage{enumitem}
\usepackage{titlesec}
\usepackage{parskip}

\hypersetup{
  colorlinks=true,
  linkcolor=black,
  urlcolor=blue
}

\titleformat{\section}{\large\bfseries}{}{0em}{}
\titleformat{\subsection}{\normalsize\bfseries}{}{0em}{}
\setlist[itemize]{leftmargin=1.5em}

\begin{document}

\begin{center}
    {\LARGE \textbf{Kaze Technical Brief}}\\[0.5em]
    {\large Architecture and Platform Overview}\\[0.75em]
    April 2026
\end{center}

\section{Purpose}

This document gives future developers and technical stakeholders a high-level understanding of Kaze before they read the code.

\section{Current System Shape}

Kaze is currently structured as a Kotlin Multiplatform product with:
\begin{itemize}
    \item a shared domain layer
    \item a Compose Multiplatform app
    \item an initial Ktor backend
    \item platform-specific targets for Android and iPhone
\end{itemize}

\section{Main Product Domains}

Current and emerging domains include:
\begin{itemize}
    \item guest experience
    \item service requests
    \item event schedules
    \item explore/discovery
    \item venue maps and wayfinding
    \item digital access pass
    \item venue reservations
    \item payments
    \item add-on services
\end{itemize}

\section{Key Architectural Direction}

Kaze is moving toward a reusable venue platform rather than a hotel-only app.

The most important bounded domains are:
\begin{itemize}
    \item \textbf{Venue Maps}: reusable spatial data for hotels and non-hotel venues
    \item \textbf{Venue Reservations}: booking conference rooms, wedding venues, and related spaces
    \item \textbf{Payments}: local payment orchestration for Rwanda
    \item \textbf{Access / Kaze Pass}: pass-based entitlements tied to venue rules
    \item \textbf{Add-On Services}: event styling, cleaning, insurance, media, and related services
\end{itemize}

\section{Venue Maps Direction}

The venue map domain should support:
\begin{itemize}
    \item sites
    \item structures/buildings
    \item levels/floors
    \item spaces
    \item nodes and paths
    \item access rules
    \item fixed objects
    \item movable layout objects
\end{itemize}

This is important because the same model can support:
\begin{itemize}
    \item hotel rooms
    \item conference halls
    \item wedding spaces
    \item stadium sections
    \item government/public buildings
\end{itemize}

\section{Layout Planning Terminology}

Two concepts should stay distinct:

\subsection{Layout Planning}

This means spatial arrangement:
\begin{itemize}
    \item chairs
    \item tables
    \item aisles
    \item seating sections
    \item fixed versus movable objects
\end{itemize}

\subsection{Event Styling}

This means visual or service decoration:
\begin{itemize}
    \item flowers
    \item lights
    \item backdrops
    \item wedding decor
\end{itemize}

Layout planning belongs near the map/space model. Event styling belongs nearer to venue commerce and add-on services.

\section{Rendering Direction}

The current UI already uses a canvas-oriented approach for venue maps. That direction remains valid for future venue layout planning.

Expected rendering characteristics:
\begin{itemize}
    \item most geometry is static
    \item spaces, walls, and fixed objects do not move continuously
    \item layout changes happen when switching from one configuration to another
    \item lightweight transitions are enough for the first versions
\end{itemize}

This is conceptually similar to scene-based or painter-driven spatial drawing: lines, rectangles, zones, symbols, and placed objects.

\section{Service Boundary Direction}

The venue-map capability is a likely candidate for a dedicated service and database because:
\begin{itemize}
    \item it can serve multiple apps
    \item not every app needs to load map concerns
    \item it can outgrow the hotel-only product context
\end{itemize}

Potential long-term services:
\begin{itemize}
    \item venue maps service
    \item reservations service
    \item payments/orchestration service
    \item guest experience app/backend
\end{itemize}

\section{Payments Direction}

Local payment support is strategically important. Early payment design should consider:
\begin{itemize}
    \item MTN MoMo
    \item Airtel Money
    \item BK / RSwitch flows
    \item reservation deposits
    \item full payment
    \item refunds and status synchronization
\end{itemize}

\section{Access Direction}

Kaze Pass should become more than a visual card. It should support entitlement logic such as:
\begin{itemize}
    \item reservation-based access
    \item payment-based access
    \item event-only access
    \item VIP or restricted zone entry
\end{itemize}

\section{Recommended Reading Order}

Before reading the code, future developers should review:
\begin{enumerate}
    \item product docs
    \item venue platform strategy
    \item PlantUML diagrams
    \item shared domain models
    \item backend boundaries
    \item UI architecture
\end{enumerate}

\section{References}

Useful project documents:
\begin{itemize}
    \item \texttt{docs/product-overview.md}
    \item \texttt{docs/venue-platform-strategy.md}
    \item \texttt{docs/map-access-control.md}
    \item \texttt{docs/diagrams/}
    \item \texttt{ROADMAP.md}
\end{itemize}

\end{document}
